% Generated draft: requires mathematical and source review.
\documentclass[11pt]{article}
\usepackage[margin=25mm]{geometry}
\usepackage{fontspec}
\setmainfont{Latin Modern Roman}
\usepackage{amsmath,amssymb,mathtools,unicode-math}
\setmathfont{Latin Modern Math}
\usepackage{enumitem,needspace,xurl,hyperref}
\setlength{\parindent}{0pt}
\setlength{\parskip}{5pt}
\setlength{\emergencystretch}{3em}
\newcommand{\N}{\mathbb{N}}
\newcommand{\Z}{\mathbb{Z}}
\newcommand{\Q}{\mathbb{Q}}
\newcommand{\R}{\mathbb{R}}
\newcommand{\C}{\mathbb{C}}
\newcommand{\F}{\mathbb{F}}
\newcommand{\A}{\mathbb{A}}
\newcommand{\PP}{\mathbb{P}}
\newcommand{\E}{\mathbb{E}}
\newcommand{\eps}{\varepsilon}
\newcommand{\dd}{\,\mathrm{d}}
\begin{document}
\section*{UnsolvedMath ID: 20000311}
\textbf{Source problem code:} catalogue entry
\subsection{Definitions and mathematical statement}
This record accidentally merges Problems 8, 9, and 10. The source-verification report specifies varieties over $\C$. For a prime divisor $E$ over a smooth variety and an ideal pair $(X,\mathfrak a^c)$, write $a_E(X,\mathfrak a^c)=A_X(E)-c\operatorname{ord}_E(\mathfrak a)$, with $A_X(E)$ the log discrepancy. Log canonicity means these numbers are nonnegative. For normal $\Q$-Gorenstein varieties use the corresponding discrepancy definition.

\begin{enumerate}
\item Problem 8 (Schwede). Let $Y$ be smooth projective and $X\subset Y$ irreducible, normal, $\Q$-Gorenstein, and of codimension $r$, with ideal sheaf $\mathcal I$. The source recalls $(Y,\mathcal I^r)$ log canonical if and only if $X$ is log canonical in the locally complete-intersection case, and says the ordinary-power equivalence fails in general. It asks whether
\[
 (Y,\mathcal I^{(r)})\text{ is log canonical}
 \quad\Longleftrightarrow\quad X\text{ is log canonical},
\]
where $\mathcal I^{(r)}$ is the symbolic power, obtained by taking the ordinary $r$th power at the generic point of $X$ and contracting. The ideal pair on the left has exponent one; symbolic power is not replaced by ordinary power.
\item Problem 9 (Musta\c t\u a). Find an adjunction formula for multiplier ideals under restriction to subvarieties not defined by a regular sequence. The relevant pair/ideal, restriction hypotheses, and desired correction terms are unspecified. The report says corrected formulas using adjoint and l.c.i.-defect ideals are available, without stating them.
\item Problem 10 (de Fernex). Let $X$ be smooth and $B\subsetneq X$ closed. Assume a prime divisor $E$ over $X$ has center $P$ and $a_E(X,cB)\leq0$ for some $c>0$. For every integer $1\leq e<\dim X$, ask for a smooth codimension-$e$ subvariety $Y$ containing $P$ and a divisor $F$ over $Y$ with center exactly $P$ such that
\[
 a_F(Y,cB|_Y-eP)\leq0.
\]
The damaged source condition ``$Y*B$'' is read by the research report as $Y\not\subseteq B$; its promised repaired transcription is not included. A valuative interpretation of the negative term would add $e\operatorname{ord}_F(\mathcal I_{P/Y})$ to $A_Y(F)-c\operatorname{ord}_F(\mathcal I_B\mathcal O_Y)$; the precise pair convention needs confirmation, especially for positive-codimension centers.
\end{enumerate}

The report says it did not find the exact symbolic-power equivalence settled, reports a corrected affirmative framework for Problem 9, and restricts known Problem 10 results to homogeneity/tangency hypotheses. Its candidate counterexample to a literal positive-dimensional-center reading is $X=\A^3$, $P=B=\{x=y=0\}$, and $c=e=2$, with $E$ the ordinary blow-up divisor. Separately, for a smooth codimension-$q$ center and that ordinary blow-up divisor, it claims that every $1\leq e<q$ admits, locally near the generic point of $P$, an order-preserving smooth slice satisfying
\[
 a_F(Y,cB|_Y-eP)=a_E(X,cB).
\]
Here $F$ is exceptional over $Y$ if $q-e\geq2$, and $F=P$ as a divisor on $Y$ if $q-e=1$. This local special-divisor identity is not the unrestricted global assertion for arbitrary $E$. All status, counterexample, and slice claims remain unreviewed.
\subsection{Short English statement}
Keep the three merged questions separate: symbolic-power inversion of adjunction, multiplier-ideal adjunction beyond regular embeddings, and preservation of a discrepancy on smooth slices. The report proposes a quantifier obstruction and a restricted local slicing identity, while giving only qualified status assessments for the other questions.
\subsection{Sources}
\begin{itemize}
\item Local UnsolvedMath registry, permanent ID $20000311$, source code \texttt{catalogue entry}; created and updated 2026-08-14. Status \texttt{partially\_solved}; classification \texttt{PARTIAL-PROGRESS}; difficulty L3; proposer field null.
\item AIM Problem Lists, Algebraic geometry, workshop \emph{Numerical invariants of singularities and higher-dimensional algebraic varieties}, source item 8 accidentally including 9 and 10, source type PDF, \url{https://aimath.org/WWN/singularvariety/problems_final.pdf}. Canonical location \texttt{aim-algebraic-geometry-notes.json}, \texttt{notes[310]}, tag \texttt{problem}; statement status \texttt{reconstructed\_unverified}; clean statement null. The verification reports the complex ground field, merged extraction, and terminal ``1'' as a page number. Original remarks and literature fields empty.
\item Source attributions: Schwede (8), Mircea Musta\c t\u a (9), de Fernex (10). The report mentions adjoint/l.c.i.-defect corrections and de Fernex homogeneity/tangency results, but supplies no exact references or hypotheses.
\item Supplied research attempt 1: result \texttt{partial\_result}, run status \texttt{partially\_solved}, rigor label \texttt{proved}, low novelty confidence. Candidate contribution \texttt{proposition}: the positive-dimensional-center counterexample and ordinary-blow-up local slice equality. Explicit caveat: machine-generated research aid, not peer reviewed; independently verify before citation.
\item Dataset attribution: Ulam AI / UnsolvedMath Contributors, CC BY 4.0; original problem attributions remain to the supplied AIM workshop collection.
\end{itemize}
\end{document}
